Since no vertex sends a positive contribution, we may keep the actual
contribution $h_x-b_x$ of each vertex of signature size one in
\eqref{eq:multi-loss-budget}. Put $\delta=9-K$ and
let $a$ count the vertices of signature size one. For an outside
vertex $x$, write $e_x$ for its marked degree outside the core under
consideration, counted with multiplicity. A vertex with $k_x=1$ and a marked
neighbor in $B$ or a marked neighbor of signature size two has loss
\[
 \ell_x=-\frac{1}{2}+2(b_x-3)+(13-h_x-b_x).
\]
Indeed $b_x\ge3$ by the comparison preceding
\eqref{eq:multi-rough-discharge}. For a marked edge from $x$ to
$U_i$, write $\{j\}$ for its signature, reflecting so that $j=i+2$.
That comparison gives at least $a_j-s_i+1=3+s_{j+2}$ blue neighbors
in $U_j$. Consequently, if $b_x=3$, all three blue neighbors lie in
$U_j$, and the other possible part $U_{j+2}$ for a marked neighbor
has $s_{j+2}=0$. It cannot receive a marked edge from $x$. A vertex
of signature size one with no such marked neighbor can send marked
edges only to the unique vertex $z$ of signature size at least three.
In this case $h_x\le Q_x-q=b_x-q$, and its loss is at least
$\frac{13}{2}+q$.

Enlarge the core by the vertices of signature size two with $b_x=1$.
The two signature positions must be nonadjacent: at adjacent positions
both anchor pairs would be blue. One of the two anchor pairs is red.
Comparison with the marked edges of the other anchor star forces
complete adjacency to its part. If a vertex in the first part were
missed, at most the remaining degree at that vertex could be used
for marked neighbors in the second part. By
\eqref{eq:multi-core-identities}, this would leave at least two blue
neighbors there. Thus these vertices enter the common neighboring part in the enlargement.
Let $n$ count the vertices of signature size two left outside.
If $b_y=2$ and $y$ has a marked neighbor in $B$, then for $L\le2$
the two signature positions cannot be adjacent. To see this, put them
at $0,1$. Both anchor pairs are blue. A marked edge into $U_2$ or
$U_4$ forces at least three blue neighbors, by the opposite anchor
star. For a marked edge $yu$ with $u\in U_3$, at least
$a_i-(s_3-1)-1\ge2$ vertices of each of $U_i$, $i=0,1$, are
adjacent to neither $y$ nor $u$. All pairs between the two omitted
sets must be blue, giving $L\ge4$.

We record the cover bounds needed for the enlarged core. Suppose
$k$ vertices have entered the five parts. Each original vertex has
at most $L+k$ blue neighbors in the two neighboring enlarged parts.
For each new vertex the bound is $k$, since its old blue degree was
one. Let $X$ be independent in the marked graph on $B'$. If the
occupied parts include an adjacent pair and form a connected set of
indices, the spanning tree argument for \eqref{eq:multi-small-cover}
gives $|X|\le L'+1$. Otherwise, if there is an adjacent occupied
pair, its two parts have $x,y$ vertices of $X$ and the remaining
occupied part is isolated. We have $x,y\le L+k$, $xy\le L'$.
If no two occupied parts are adjacent, then $|X|\le13$.
The five bounds on $|X|$ used below follow by substituting
the values of $R',k,L,L'$ in these inequalities. They are listed in
\textbf{Table~\ref{tab:thirteen-cover-check}} in Appendix~\ref{app:multi-structure}.
In the fourth row an independent set of size 13 consists of two entire
nonadjacent parts. This also
holds if $L'\le1$ and $\max_i a_i'\le10$, since an independent set
using an adjacent pair then contains at most $10+1+1$ vertices.
The same argument applies to $B$ with $L\le1$ and a vertex of
signature size three. In each of these situations, if $xy$ is marked
outside and $k_x+k_y\ge4$, then
\begin{equation}\label{eq:d13-improved-cover}
 e_x+e_y\le 26-(|B'|-13)-1.
\end{equation}
Indeed equality in the ordinary cover bound would force the two core
neighborhoods to be disjoint and their complement to be an independent
set of size 13. Their union would then be contained in three parts,
whereas disjoint neighborhoods contain $k_x+k_y$ distinct anchors.

Suppose $k_x=1$, $b_x=3$, and $x$ has a marked neighbor in $B$.
Rotate and reflect so that its signature is $\{3\}$ and its only part
for marked neighbors is $U_0$. It follows that $s_1=0$. If $xy$ is marked and
$k_y=2$, its two positions must be $\{0,2\}$: the other possibilities
$\{0,1\},\{1,4\}$ giving a cover with $3$ use position 1.
The marked stars from the two omitted anchors require $x,y$ together
to see $U_0,U_2,U_3$. Hence $y$ sees all of $U_2$, and misses a
vertex of $U_0$ because $a_0+a_2=13$ and $xy$ uses one degree.
Comparison with the marked spoke at the missed vertex bounds the
number of red neighbors in $U_2$ by $s_0$. Writing $Q_y(U)=
|N_{G_{\rm s}}(y)\cap U|$, we obtain
\begin{equation}\label{eq:d13-two-positions}
 |\,\{u\in U_2:yu\text{ is blue}\}\,|\ge2+s_4,\qquad
 Q_y(U_3)\ge a_3-3=s_0-1.
\end{equation}
If $b_y=2$ then $s_4=0$ and both blue neighbors lie in $U_2$.
For $s_0\ge2$ there is a red neighbor in $U_3$. Comparison of this
edge with the marked star from $c_1$ to $U_0$ gives
$Q_y(U_0)\ge a_0-s_3+1=3+s_2$. Summing the three part counts gives
\begin{equation}\label{eq:d13-outside-degree}
 b_y=2,\ s_0\ge2\quad\Longrightarrow\quad e_y\le9-2s_0-s_2 .
\end{equation}
If instead $xz$ is marked with $k_z=3$, the possible signatures
of $z$ are $\{0,2,3\}$ and $\{0,2,4\}$. Both omit position 1.
Panels (a) and (b) of \textbf{Figure~\ref{fig:degree-thirteen-signatures}}
show the two possibilities.
The star from $c_1$ forces $z$ to see every vertex of $U_2$, as
well as two anchors outside that part. Thus $Q_z\ge a_2+2$.
In particular the signature of such a $z$ is itself a cover of the cycle.

Suppose first that all outside vertices have at most two positions.
There is no vertex of signature size zero. The losses of the $a$
vertices with one position and the $\frac{9-\delta-a}{2}$ vertices with
two positions are at least $-\frac{1}{2}$ and 2 per vertex, respectively.
Subtracting these amounts from \eqref{eq:multi-loss-budget} gives
the following upper bound on $2L$ plus all excess losses:
\[
 T=\frac{3}{2}+\frac{3a}{2}-\frac{11\delta}{2}.
\]
Each remaining vertex with two positions costs at least
$2(b_x-1)\ge2$ beyond its minimum, and a vertex with one position
costs at least $2(b_x-3)$. When $n=0$ there is no outside marked
edge. If $a\ge2$ then $L_{\rm cut}\ge6$, and
\eqref{eq:multi-extra-capacity} contradicts $\mu\le\frac{21}{4}$. If $a=1$
with $\delta=0$, the base count for $R'=1$ applies.
For $\delta=2$ the budget forces $a=7$, $n=0$.
For $\delta=1$, $a=4,6,8$. When $a=8$ we have $n=0$. When $a=6$
we have $n\le1$ and $F_{\rm out}\le13$. If $n=1$, the cut
has at least $18+2$ blue pairs. When $a=4$, $T=2$ and $n\le1$.
If $n=1$ the other vertex with two positions was absorbed, $L\le1$,
and the old cover number is at least 15. Equation
\eqref{eq:multi-center-cut} gives $F_{\rm out}\le8$, against
$L_{\rm cut}\ge12+2$.

It remains in this case to take $\delta=0$ and $a$ odd. If $a\ge7$,
then $n\le1$, $F_{\rm out}\le13$, and $L_{\rm cut}\ge21$.
If $a=1$, then $T=3$ and $n\le1$. For $n=1$, write $x,y$ for the
vertices with one and two positions. We have $L=0$, $b_x=3$,
$h_x\ge9$, $b_y=2$, $h_y\ge10$, and the other three vertices
with two positions entered the core. The first row of \textbf{Table~\ref{tab:thirteen-cover-check}}
gives a cover of size at least 18, so $F_{\rm out}\le4$. If
$F_{\rm out}=0$, then $L'\ge3$ and $L_{\rm cut}\ge5$ suffice.
Otherwise both vertices have marked neighbors in $B'$. Their three
positions form a cover, say $\{0,2,3\}$, with remaining degree one
at each position. Compare the stars at the omitted anchors as in the $R'=3$ case
for degrees 11 and 15, now with $d=6$. If $x$ uses position 0, it gives
$\mu\ge13-2w_0$, where the number $w_0$ of blue edges of the
chosen star is at most three, one per added vertex. If $x$ uses a
leaf position and has a marked neighbor carrying the required red
spoke, it gives $\mu\ge12-F_{\rm out}$. In the other case every
marked copy from $x$ to $B'$ uses a distinct blue spoke and
$\mu\ge h_x+b_x+b_y-2F_{\rm out}\ge6$.

For $a=3$ the budget is $T=6$. If $n=1$, then $R'=5$, $L\le2$,
$L'\le6$. The second row gives a cover of size at least 16, so
$F_{\rm out}\le7$ and $L'+L_{\rm cut}-F_{\rm out}\ge2+9+2-7=6$.
If $n=2$, then $R'=7$, $L\le1$, $L'\le3$ and the third row gives
a cover of size at least 16. If the two remaining vertices with two
positions are joined by $w\ge1$ marked copies, their outside degrees
sum to $F_{\rm out}+w\le10$. The only way
$L'+L_{\rm cut}-F_{\rm out}<6$ is $F_{\rm out}=9$, $L'=1$,
$L_{\rm cut}=13$. Then $w=1$ contradicts
\eqref{eq:d13-improved-cover}. Otherwise all outside marked edges
join a vertex with one position to a vertex with two. At least one
of the latter vertices has blue degree two. If none of blue degree
two is joined by a mark to a vertex of blue degree three with one
position, all outside marks meet one vertex. Indeed if both blue
degrees are two, at most one vertex with one position has blue
degree greater than three. If one blue degree is three, the two
vertices use the entire budget and all vertices with one position
have blue degree three. Equation \eqref{eq:multi-center-cut} gives
$F_{\rm out}\le7$ in either case.

Thus let $y$ have $k_y=b_y=2$ and a marked neighbor $x$ with
$k_x=1,b_x=3$. Here $h_x\ge8$. If $x$ had no marked neighbor in
$B$, it could receive at most one marked copy from each outside
neighbor, by $q=2$ and $Q_x=3$, which is impossible. By
\eqref{eq:d13-two-positions} we may take $y$ at $\{0,2\}$ and
$s_1=s_4=0$. The absorbed vertex with two positions must use 0.
If the other remaining vertex $y'$ has a marked neighbor with one
position and blue degree three, its positions are also nonadjacent
and include 0. Hence $s_0\ge3$, and \eqref{eq:d13-outside-degree}
gives $e_y\le2$. If $b_{y'}=2$ the same bound gives $e_{y'}\le2$.
If $b_{y'}=3$, the star bound with three possible outside neighbors
gives $e_{y'}\le7$, and the cut includes an additional blue pair.
If $y'$ has no such marked neighbor and has positive outside marked
degree, both blue degrees are two. At most one vertex with one position
has blue degree greater than three, so $e_{y'}\le5$. Both vertices
have old core marks since $q=2$ excludes purely outside marks at
blue degree two, and both have nonadjacent positions. Again $s_0\ge3$.
It follows throughout that $F_{\rm out}\le8$, or that
$F_{\rm out}\le9$ and $L_{\rm cut}\ge14$.
\eqref{eq:multi-extra-capacity} now gives $\mu\ge6$.

If $a=3,n=3$, then $L=0$. Each vertex with one position has
$(b_x,h_x)=(3,10)$ and each with two has $(2,11)$. Each of the
latter has an old core mark and nonadjacent positions.
Whenever an outside mark meets a vertex with one position,
\eqref{eq:d13-two-positions} permits only positions $0,2,3$ after
rotation and reflection, so all three vertices with two positions
use 0. Each one joined by a mark to a vertex with one position
has outside degree at most two by \eqref{eq:d13-outside-degree}.
If there is just one such vertex, the other two share at most five
marks, giving $F_{\rm out}\le7$. If there are at least two, all
outside marks meet these vertices, giving $F_{\rm out}\le6$. If
there are none, the outside marks lie on three vertices. On a triangle,
summing the three degree-sum bounds gives $4F_{\rm out}\le33$.
For a star the degree bound gives $F_{\rm out}\le7$. In all cases
$F_{\rm out}\le8$ and $L_{\rm cut}\ge15$.

For $a=5$ we have $T=9$ and originally two vertices with two
positions. If $n\le1$ every outside mark meets a single such vertex.
Since $L\le4$, the old cover number is at least 14. The star bound
with five possible neighbors gives $F_{\rm out}\le10$. If $n=1$,
the absorbed vertex gives $L'\ge1$ and $L_{\rm cut}\ge17$.
Take $n=2$. Now $L\le2$, the old cover number is at least 15,
and $L_{\rm cut}\ge19$. If the two remaining vertices share a mark,
$F_{\rm out}\le10$. Otherwise the outside marks form a bipartite
graph between these two vertices and the five vertices with one
position. The outside marked degrees at the ends of each mark sum
to at most 11.

If every vertex with one position incident with an outside mark has
blue degree at least four, the budget allows at most two such vertices.
For two disjoint outside marks, their degree sums yield
$2F_{\rm out}\le22$. Otherwise the outside marks form a star with
at most two neighbors and $F_{\rm out}\le7$. Now suppose a vertex
$x$ with $k_x=1,b_x=3$ meets an outside mark. It must have an old
core mark since $h_x\ge5$ by the budget, whereas $Q_x=3$ would
allow at most two outside marked copies in the absence of a core mark.
If a neighbor $y$ with $k_y=2$ has $b_y=2$, take its positions to
be $\{0,2\}$ with $s_1=s_4=0$ as in \eqref{eq:d13-two-positions}.
The other vertex $y'$ with two positions, if these are adjacent, must
use $\{2,3\}$ and have $b_{y'}\ge3$. It can meet outside marks only
at vertices with blue degree greater than three, so it has at most
one such neighbor, and $e_{y'}\le5$. Also $e_y\le9$, since $y$
sees $U_2$ and $c_0$. When $e_{y'}>0$ the cut has at least 21
blue pairs. When $e_{y'}=0$, $F_{\rm out}\le9$. If instead the positions
of $y'$ are nonadjacent, both $y$ and $y'$ use 0, giving $e_y\le4$
by \eqref{eq:d13-outside-degree}. Since $e_{y'}\le10$, cut size
20 suffices. At size 19 both blue degrees are two and all five vertices
with one position have blue degree three, so $e_{y'}\le4$ as well.

If vertices with one position and blue degree three meet outside marks
only at the vertex of blue degree at least three, the other vertex $y'$
of signature size two has blue degree two. It has at most one possible
marked neighbor outside the core, so $e_{y'}\le5$. When its degree is positive the cut has at least 21 blue
pairs and $F_{\rm out}\le15$. When its degree is zero, $F_{\rm out}\le10$.
The only case left has both blue degrees at least three. The budget
forces both to be three, $L=0$, and blue degree three for all five
vertices with one position. If $F_{\rm out}\ge16$, one vertex $y$
with two positions has $e_y\ge8$ and the other also has a marked
outside neighbor. From \eqref{eq:d13-two-positions} applied at $y$,
we have $s_1=0$, $s_4\le1$ and $Q_y\le5$. The whole part $U_2$,
the vertex $c_0$, and the $s_0-1$ or more neighbors in $U_3$ give
$Q_y\ge2+2s_0+s_4$, so $s_0=1$. The other vertex with two positions
must have nonadjacent positions if it meets a vertex with one
position and blue degree three, by the same signature comparison.
It cannot use 0 or 1 and hence uses $\{2,4\}$. The comparison then
requires $s_3=0$, whereas $y$ already has a neighbor of signature
$\{3\}$. Thus $F_{\rm out}\le15$ and the 21 blue pairs suffice.

We next consider the case of a unique outside vertex $z$ of signature
size at least three. Write $n_2$ for the number with two positions
before enlargement. From $k_z+a+2n_2=9-\delta$ and
\eqref{eq:multi-loss-budget}, the amount left for $2L$ and the
losses above their respective lower bounds is at most
\[
 \frac{21}{2}-\frac{13\delta}{2}
 -13\left(\frac{k_z}{2}-1\right)+\frac{a}{2}-2n_2 .
\]
The nonnegative possibilities are
\[
\begin{array}{rrrrr}
 \delta& k_z&a&n_2&\text{remaining amount}\\
 1&3&5&0&0\\
 0&4&5&0&0\\
 0&3&2&2&1\\
 0&3&4&1&4\\
 0&3&6&0&7
\end{array}
\]
Every vertex of signature size one has an old core mark or a marked
neighbor of signature size two, since otherwise its loss exceeds its
lower bound by at least $7+q\ge8$. In particular $b_x\ge3$ for
every such vertex. When $n_2=0$, every outside mark meets $z$.
An outside neighbor with signature zero has only blue neighbors in
the core and at least $e_z+q$ of them by the cover bound. These
already pay for all $F_{\rm out}=e_z$ marks and there remain the
blue cut pairs of the vertices with one position. If there is
no such neighbor, then $a=5$, $L=0$ gives a cover of size at least
15 and hence $F_{\rm out}\le9$, while $a=6$ gives $e_z\le10$ since
$z$ sees at least three anchors. The cut has at least $3a$ blue pairs.

For $a=2$, both vertices with two positions have blue degree one
and enter the core. Thus $R'=5$, $L=0$, $L'\le3$. The two vertices
with one position have blue degree three and combined marked degree at
least 19. Both have old core marks, since otherwise either could have
at most three outside marks, against marked degree at least nine.
The fourth row of \textbf{Table~\ref{tab:thirteen-cover-check}} gives $\tau'\ge17$. A marked
neighbor of $z$ with signature zero uses at least $e_z+4$ blue
cut pairs. Together with the six supplied by the two vertices with
one position, these suffice. Otherwise the only possible neighbors
of $z$ outside are the two vertices with one position.
Equation \eqref{eq:d13-improved-cover} gives outside degree sum at
most eight for each such pair, whence $F_{\rm out}\le5$ by the star
bound. At a vertex $x$ with one position $j$, a mark to its sole
possible part $U_i'$ forces at least
$a_j'-s_i'+1-h=5-h$ blue neighbors, where $h$ counts the blue pairs
in the opposite anchor star into $U_j'$. This is the same comparison
used in the case $D=11,15$: at most $s_i'-1$ neighbors can be
supplied by the other end of the mark, and only the $h$ blue
pairs can fail to impose an adjacency. The other possible marked
part has old and new remaining degree zero. If both vertices used
the same star they would have the same signature and the same marked
part $U_i'$ with $s_i'\le1$. Since $a_i'\le8$, this part cannot
supply their required $19-F_{\rm out}\ge14$ marks. Distinct stars
have disjoint sets of blue pairs, so $L'+L_{\rm cut}\ge10$.

By \eqref{eq:multi-extra-capacity}, the only possibility left is
$F_{\rm out}=5$ and
$\frac{R'^2}{4}-\sum_i s_i's_{i+2}'=\frac{1}{4}$. Thus $\sum_i s_i's_{i+2}'=6$.
The five entries sum to five, and at least three are positive.
Five positive entries give sum of products five. With four, their
support in the graph with edges $\{i,i+2\}$ is a path with one entry
two, again giving at most five. With three, reaching six requires a
two-edge star with center weight two or three. Put its center at 0
and its leaves at 2 and 3. The two choices of weights give
$(0,1,2,0,2)$ and $(0,1,1,0,3)$ up to rotation and reflection. Both
appear in the $R=5$ group of \textbf{Figure~\ref{fig:finite-simple-slack-low}}.
The three positions of $z$ are exactly
these. Neither of the vertices with one position can use 0, since
both its possible marked parts then have positive remaining degree.
Only $z$ can use 0, so $s_0'=1$ because $R'=K'$. This contradicts
the center weight.

Suppose $a=4$ and the single vertex with two positions enters the
core. Then $R'=7$, $L\le2$, $L'\le5$, and the final row of
\textbf{Table~\ref{tab:thirteen-cover-check}} gives $\tau'\ge15$. A marked neighbor of $z$ of signature
zero again supplies at least $e_z+2$ blue cut pairs. Otherwise
$F_{\rm out}\le8$ by the star bound. Since $L'\ge1$ and
$L_{\rm cut}\ge12$, only $F_{\rm out}=8$, $L'=1$, $L_{\rm cut}=12$
can have $\mu<6$. All four vertices with one position have blue
degree three, and all have an old core mark. Indeed, the old cover
is at least 15 and in the absence of an old mark they can receive
at most two copies outside, against $h_x\ge6$ from the loss bound.
The preceding anchor comparison gives blue degree at least $4-h$
for each, so all four must use the unique missing edge in the same
star. Their sole marked part has size at most ten and remaining degree
at most one. It would have to supply at least $4\cdot10-4-8=28$
marked copies, using the total excess loss at most four. This is
impossible.

Finally let $a=4$ with $z$ of signature size three and $y$ of
signature size two still outside. There was no enlargement. Here
$\delta=0$, $b_y\ge2$, $L\le1$, and the cover number is at least
15. The improved cover bound gives $e_z+e_x\le10$ whenever $zx$
is marked outside and $k_x>0$. Each vertex with one position and
blue degree three has an old core mark, since $h_x\ge8$, whereas
without a core mark it could have at most two outside copies.
The five vertices with positive signature other than $z$ supply
at least $12+2=14$ blue cut pairs. If $z$ has a marked neighbor with
signature zero, its blue degree is at least $e_z+2$. The outside
marks not meeting $z$ form a star at $y$ with at most four neighbors.
There are at most $\lfloor\frac{4\cdot11}{5}\rfloor=8$ such copies, using
the ordinary cover bound for their endpoints. Thus $\mu\ge14+2-8$.
We may now assume there is no marked neighbor with signature zero and
hence no vertex of signature zero in the reduced host.

If $z,y$ are joined by $w\ge1$ marked copies, then
$F_{\rm out}=e_y+e_z-w\le9$. Only $L=0$, $L_{\rm cut}=14$,
$F_{\rm out}=9$ needs consideration. Thus $b_y=2$, and all four
vertices with one position have blue degree three. If $y$ has such a
marked neighbor, take its signature to be $\{0,2\}$ by
\eqref{eq:d13-two-positions}. Then $s_1=s_4=0$, $z$ uses $\{0,2,3\}$,
and the four vertices can only use 2 or 3 and mark into $U_0$.
All remaining degree of the old anchors comes from outside signatures
since $K=R$. Consequently $s_0=2$. The capacity of $U_0$ is
$s_0a_0=s_0(11-s_0)=18$, whereas these vertices need at least
$4\cdot10-2-9=29$ core marks. If $y$ has no such marked neighbor,
all outside marks meet $z$, so $e_z=9$ and each of its at most five
neighbors has outside degree at most one. This too is impossible.

Suppose $y,z$ have no marked edge between them. The outside marks
are then bipartite between $\{y,z\}$ and the four vertices with one
position, giving $F_{\rm out}=e_y+e_z$. The budget excludes $b_y>3$.
If $b_y=3$, then $L=0$ and each of the four vertices has blue degree
three and marked degree ten. We need only exclude $F_{\rm out}\ge10$.
If either $e_y$ or $e_z$ is zero, the star bound gives
$F_{\rm out}\le8$. Otherwise both have a marked neighbor with one
position and blue degree three. By the three-position comparison
following \eqref{eq:d13-outside-degree}, rotate the signature of $z$
to $\{0,2,3\}$, with $s_1=0$ or $s_4=0$. If both vanish, all four
vertices with one position use 2 or 3 and mark into $U_0$.
Vertex $y$ has positions $\{0,2\}$ or $\{0,3\}$. Thus $s_0=2$
and the capacity 18 of $U_0$ is smaller than $40-e_y-e_z$, since
$e_y\le10$ and $e_z\le9$.

Otherwise reflect to make $s_1=0$ and $s_4>0$. The four signatures
are 3 or 4, giving marks to $U_0$ or $U_2$, respectively. Vertex
$y$ must use $\{0,2\}$, and $s_0=s_2=2$. If $y$ has a marked neighbor
of signature $\{3\}$, \eqref{eq:d13-two-positions} gives $s_4\le1$,
so $s_4=1$. Its three blue neighbors all lie in $U_2$ and it has a
marked neighbor in $U_3$. Comparison with the star from $c_1$
therefore gives $Q_y(U_0)\ge3+s_2=5$. We also have $Q_y(U_2)=5$ and
$Q_y(U_3)\ge1$. Hence $e_y\le2$. For a neighbor of signature
$\{4\}$ interchange $0,2$ and $3,4$ to get the same bound.
A marked neighbor of $z$ with one position forces $z$ to see a
whole part $U_0$ or $U_2$ of size at least five, and two additional
anchors. Thus $e_z\le6$, contradicting $F_{\rm out}\ge10$.

We are left with $b_y=2$. There is at most one vertex $x_0$ with
one position and blue degree greater than three. When it exists,
$b_{x_0}=4$, $h_{x_0}=9$, $L=0$, and the other three vertices have
blue degree three and marked degree ten. Here $F_{\rm out}\ge10$
is necessary. First suppose $y$ has a marked neighbor of blue degree
three with one position. By \eqref{eq:d13-two-positions}, $s_1=s_4=0$
and $y,z$ have
signatures $\{0,2\}$ and $\{0,2,3\}$. Every vertex with one position
and blue degree three marks into $U_0$. If $x_0$ does not exist,
$s_0=2$, $e_y\le3$ by \eqref{eq:d13-outside-degree}, and
$e_z\le7$ because it sees a whole part of size at least four as
well as two other anchors whenever it meets an outside mark.
The required $40-2-10=28$ or more core marks exceed the capacity
18 of $U_0$. If $x_0$ exists and $z$ has no other marked neighbor,
$e_z\le5$ by $e_z+e_{x_0}\le10$, so $F_{\rm out}\le8$.
Otherwise $e_z\le7$. If $x_0$ does not use 0, then $s_0=2$ and
the other three vertices need at least $30-10=20$ core marks in
$U_0$. If $x_0$ uses 0 then $s_0=3$ and
\eqref{eq:d13-outside-degree} gives $e_y\le1$. Either possibility
is excluded.

It remains that $y$ has no marked neighbor of blue degree three
with one position. If there is no $x_0$, every outside mark meets
$z$ and $F_{\rm out}\le8$. Otherwise only $yx_0$ can contribute
to $e_y$, and its multiplicity $w$ satisfies $2w\le11$. If $zx_0$
is marked, then $e_z+w\le9$ by the improved bound and the copy on
$zx_0$. If it is unmarked, $z$ can meet only the other three vertices,
so $e_z\le\lfloor\frac{3\cdot10}{4}\rfloor=7$. To have $F_{\rm out}\ge10$
requires $w\ge3$. Now $h_y=11$ and $h_y-w\ge6$ copies go to $B$.
The positions of $y$ are nonadjacent because $L=0$ and $b_y=2$.
By the comparison following \eqref{eq:d13-outside-degree}, rotate
and reflect so that $z$ has signature $\{0,2,3\}$ with
$s_1=0$ or $s_4=0$. If both vanish, all three vertices of blue
degree three send their old core marks to $U_0$. The required
capacity is at least $30-e_z\ge23$, hence $s_0(11-s_0)\ge23$
forces $s_0\ge3$. Only $y,z,x_0$ can use 0. Thus both $y$ and
$x_0$ use it, so their signatures together have at most two positions
and cannot cover the cycle.

If just one of $s_1,s_4$ is zero, take $s_1=0$ by reflection.
The three vertices of blue degree three use 3 or 4 and mark into
$U_0$ or $U_2$, respectively. We have $F_{\rm out}=e_z+w\le12$.
There are just three possible nonadjacent pairs of positions for $y$.
Together with the position of $x_0$, they give the four cases in
\textbf{Figure~\ref{fig:degree-thirteen-signatures}}(c)--(f).

\begin{figure}[htbp]
\centering
% The last argument is the index of the singleton neighbor, or -1.
\providecommand{\cycleanchorneighbors}[4]{%
  \begin{tikzpicture}[x=1cm,y=1cm,baseline=(current bounding box.center)]
    \path[use as bounding box] (-1.38,-1.38) rectangle (1.38,1.4);
    \foreach \i in {0,...,4}
      \coordinate (c\i) at ({90-72*\i}:.83cm);
    \draw[marked edge] (c0)--(c1)--(c2)--(c3)--(c4)--cycle;
    \foreach \i in {#2} \draw[black!75,line width=.6pt] (0,0)--(c\i);
    \node[fill=white,inner sep=1pt,font=\small] at (0,0) {$#1$};
    \ifnum#4>-1\relax
      \coordinate (outside) at ({90-72*(#4)}:1.27cm);
      \draw[black!75,line width=.6pt] (c#4)--(outside);
      \node[fill=white,inner sep=1pt,font=\small] at (outside) {$#3$};
    \fi
    \foreach \i in {0,...,4}
      \node[fill=white,inner sep=1pt,font=\scriptsize] at (c\i) {$\i$};
  \end{tikzpicture}}

\begin{tabular}{ccc}
\cycleanchorneighbors{z}{0,2,3}{x}{3} &
\cycleanchorneighbors{z}{0,2,4}{x}{3} &
\cycleanchorneighbors{y}{0,2}{x_0}{3} \\
(a) & (b) & (c) \\[2pt]
\cycleanchorneighbors{y}{0,2}{x_0}{4} &
\cycleanchorneighbors{y}{0,3}{x_0}{2} &
\cycleanchorneighbors{y}{2,4}{x_0}{0} \\
(d) & (e) & (f)
\end{tabular}
\caption{Anchor neighbors for $D=13$ and $R=9$, with $c_i$ labeled $i$.
In (a) and (b), $xz$ is marked, $k_z=b_x=3$, and $x$ has signature
$\{3\}$ and a marked neighbor in $U_0$. In (c)--(f), $yx_0$
is marked, $s_1=0<s_4$, and $z$ has signature $\{0,2,3\}$. Only edges
incident with an anchor are drawn.}
\label{fig:degree-thirteen-signatures}
\end{figure}


If $y$ uses $\{0,2\}$, $x_0$ must use 3 or 4. Here $s_0=s_2=2$.
Suppose $x_0$ uses 3. Its only possible core marks go to $U_0$:
a mark into $U_1$ is excluded by $s_1=0$. It has at least three
blue neighbors in $U_3$ and at most one in other parts. Since the
marked edge $yx_0$ must meet or have adjacent ends to each edge of
the omitted anchor stars, $y$ misses at most one point of $U_2$.
It cannot see all of $U_0$, for then $Q_y\ge a_0+a_2-1=12$,
contrary to $w\ge3$. Comparison with a missed point in $U_0$
bounds its distinct marked neighbors in $U_2$ by $s_0=2$, and it
has only two blue neighbors. Thus $a_2\le5$, forcing $s_4=1$.
Two of the three vertices of blue degree three use 3 and one uses 4.
Together the two vertices using 3 and $x_0$ need at least
$29-F_{\rm out}\ge17$ marks into $U_0$, whose capacity is
$s_0a_0=2(4+4)=16$. If $x_0$ uses 4, the reflected comparison
gives $s_3\le1$. Position 3 already occurs at $z$. Thus all three
vertices use 4 and at least $39-F_{\rm out}$ marks must go to $U_2$,
again of capacity 16.

If $y$ uses $\{0,3\}$, then $x_0$ uses 2 to cover the cycle with
$y$. Since $s_0=2$, a mark from $x_0$ to $U_4$ would force
$b_{x_0}\ge3+s_0=5$. It can mark only into $U_0$, and the same
star comparison gives $s_4\le1$. Thus $s_4=1$, $s_2=2$, $s_3=4$,
and two of the three vertices of blue degree three use 3. Again
at least $29-F_{\rm out}$ marks must enter $U_0$ of capacity 16.
If $y$ uses $\{2,4\}$, $x_0$ uses 0. Now $s_2=2$ forces its
core marks into $U_2$ and $s_3\le1$. All three vertices of blue
degree three use 4 and $s_0=2$, $s_4=4$. At least $39-F_{\rm out}$
marks must enter $U_2$ of capacity 16. These contradictions exhaust
the possibilities for $(D,R)=(13,9)$.
